\documentclass{report}
\usepackage{amsfonts, amsmath}

\newcommand\N{{\mathbb N}}
\newcommand\Q{{\mathbb Q}}
\newcommand\R{{\mathbb R}}
\newcommand\Z{{\mathbb Z}}

\pagestyle{empty}
\begin{document}
\large
\centerline{\Huge \bf Analisi Matematica II}
\vskip.2cm
\centerline{\Huge  Primo Appello (19-06-2002)}
\renewcommand{\thefootnote}{ } 
\footnote{\sl Avete 3 ore di tempo. Ogni esercizio
vale sei punti. La sufficienza si ottiene con un punteggio $\geq$
18. {\bf Solo le risposte chiaramente giustificate saranno prese in
considerazione.} {\scshape Elaborati illeggibili o confusi verranno ignorati.}
{\sffamily La
sintesi sar\`a particolarmente apprezzata.}}
\renewcommand{\thefootnote}{\arabic{footnote}} 

\vskip1cm
\begin{enumerate}
\item Si consideri la funzione
\[
f(x):=\int_0^1e^{-xy}\sqrt{1-y^2}dy.
\]
Si dimostri che \`e derivabile e che la derivata \`e sempre negativa.
\item Si calcoli
\[
\int_{-\frac\pi 2}^{\frac \pi 2} e^{\sin x} dx
\]
con una precisione superiore a 1/2.
\item Si dimostri che la serie 
\[
\sum_{n=2}^\infty \frac 1{n(\ln n)^2}
\]
\`e convergente.
\item Si risolva il seguente problema di Cauchy
\begin{align*}
&\frac {d^2x}{dt^2}=2x- \frac{dx}{dt}\\
&x(0)=0\;;\quad
\frac{dx}{dt}(0)=1.
\end{align*}
\item  Si risolva il seguente problema di Cauchy
\[
\frac {dx}{dt}=x+\cos t\;;\quad
x(0)=0.
\]
\item Si calcoli il raggio di convergenza della serie di potenze
\[
\sum_{n=0}^\infty \cos n\; x^n.
\]
\end{enumerate}
\end{document}
